\src{106}
\begin{enumerate}
\item[(i)] Một cặp ánh xạ
\[
  p_1,p_2:R\rightrightarrows X
\]
được gọi là \emph{ánh xạ nguồn} và \emph{ánh xạ đích},

\item[(ii)] Một ánh xạ
\[
  \pi:(R,p_2)\times_X(R,p_1)\longrightarrow R
\]
được gọi là \emph{ánh xạ hợp thành}.
\end{enumerate}
Các dữ liệu này phải thỏa mãn những tiên đề quen thuộc, mà ta sẽ không nhắc
lại ở đây, và có thể được diễn đạt bằng tính giao hoán của một số biểu đồ,
cùng với sự tồn tại của một ánh xạ $D:X\to R$ (nhất thiết duy nhất) làm cho
hai biểu đồ khác giao hoán. Ánh xạ $D$ gửi một vật tới mũi tên đồng nhất tương ứng
và thỏa mãn
\[
  p_1D=p_2D=\operatorname{id}_X.
\]
Việc nói phạm trù là một groupoid khi đó có nghĩa là tồn tại một ánh xạ
\[
  s:R\longrightarrow R
\]
(nhất thiết duy nhất), gọi là \emph{phép lấy nghịch đảo} của $R$, biến mỗi
mũi tên thành một mũi tên nghịch đảo. Điều này có thể được diễn đạt bằng tính
giao hoán của bốn biểu đồ khác được tạo từ $s$, $\Delta$ và các dữ liệu trước
đó; hai biểu đồ đầu là
\[
  p_1s=p_2,\qquad p_2s=p_1.
\]

Nhắc lại các điểm này, các định nghĩa tổng quát trong ([2], A, \S~1) đặc biệt
cho thấy thế nào là một cấu trúc $\mathcal C$-\emph{phạm trù}, tương ứng một
cấu trúc $\mathcal C$-\emph{groupoid}, trên một cặp vật $X,R$ của một phạm trù
tùy ý $\mathcal C$: theo định nghĩa, đó là dữ liệu, với mỗi vật $T$ của
$\mathcal C$, của một cấu trúc phạm trù, tương ứng một cấu trúc groupoid, theo
nghĩa tập hợp, có tập các vật là $X(T)$ và tập các mũi tên là $R(T)$, các cấu
trúc này \guillemotleft biến thiên theo hàm tử\guillemotright{} khi $T$ biến
thiên. Do đó, điều này xác định hai cấu xạ
\[
  p_1,p_2:R\rightrightarrows X
\]
được gọi là \emph{cấu xạ nguồn} và \emph{cấu xạ đích}, và khi tích thớ cần
thiết tồn tại, một cấu xạ
